\chapter{Conclusion}
\label{chap:conclusion}

This thesis set out to explore whether a pull-based orchestration model could offer a simpler, lighter alternative to platforms like Kubernetes and Nomad, without giving up the fault tolerance that distributed workloads demand. Through the Design Science Research methodology, the work moved from problem analysis to a working prototype and finally to empirical evaluation on real infrastructure. The results, while limited in scale, suggest that the approach has merit for the right kind of workload.

\section{Summary of Findings}

Chapter~\ref{chap:intro} posed two research questions. The first asked whether a pull-based system can meaningfully reduce operational overhead while still recovering from failures quickly. The evaluation in Chapter~\ref{chap:eval} gives a reasonably clear answer: the control plane runs on 0.5 vCPU and 1 GB of RAM, significantly less than a highly available Kubernetes setup (see Table~\ref{tab:resource-comparison}), yet it recovers from node failures in under 48 seconds on average, well within the 90-second target established in Chapter~\ref{chap:met}. Across 700 workload placements in five test scenarios, not a single lease conflict or duplicate acquisition occurred, which speaks to the correctness of the coordination mechanism.

The second research question concerned the architectural trade-offs between pull and push models. What emerges from the analysis is a consistent pattern of trading sophistication for simplicity. Because the control plane is stateless and agents pull work on their own schedule, there is no need for consensus protocols or cluster state replication, but this also means the scheduler is limited to first-come-first-served placement. It cannot make globally optimal decisions the way systems like Firmament~\cite{gog2016firmament} or TetriSched~\cite{tumanov2016tetrisched} can. The 30-second lease timeout keeps failure detection simple, but it also sets a hard floor on how quickly the system can react to outages. As discussed in Chapter~\ref{chap:met}, the design favors availability over strong consistency, which is a reasonable choice for short-lived batch workloads where a brief scheduling delay is tolerable. One finding worth noting is that scheduling multiple workloads on a single node is roughly 76\% slower than distributing them. The agent has to handle container creation sequentially while also refreshing leases and reconciling state, so the system clearly works better when workloads are spread across nodes.

\section{Strengths, Limitations, and Applicability}

The main advantage of the proposed system is its simplicity. The pull model sidesteps several failure modes that plague push-based architectures: broken deliveries after network partitions, duplicate scheduling from stale controller state, and the operational burden of maintaining a complex control plane. The lease mechanism itself is well-understood and straightforward, drawing on established distributed systems principles~\cite{gray1989leases}. Security was also kept in scope: all communication between components uses mTLS, so the system does not cut corners on that front despite its minimalism.

These properties make the system a reasonable fit for a handful of practical scenarios. Short-lived batch jobs and CI/CD pipeline workloads, where the time spent starting a container is small compared to the time it runs. Edge and IoT deployments, where the hardware simply cannot accommodate the resource demands of a full Kubernetes stack. Small teams and organizations that need orchestration but cannot justify hiring dedicated platform engineers. The business scenario from Chapter~\ref{chap:lr}, a B2B authentication service recovering from infrastructure failures, is itself a concrete example of where this approach would be appropriate.

That said, the limitations are real and should not be understated. The evaluation covered only two agents with at most two concurrent workloads per node, and it would be premature to assume the same behavior at larger scales. Database contention and polling overhead may become bottlenecks that were not visible here. The system lacks features that many production environments consider essential: rolling updates, autoscaling, affinity rules, and sophisticated placement strategies. Test workloads used minimal container images, so the scheduling times measured may not reflect how real applications behave. All testing was done within a single cloud region, meaning the effects of network latency in geographically distributed setups remain unknown.

\section{Future Work}

Several directions deserve attention. The most immediate is testing at larger scale, with dozens of agents and hundreds of workloads, to see whether the linear scaling properties hold under real production load. Long-running tests over days or weeks would help surface memory leaks and resource exhaustion issues that short trials cannot catch. Making the lease timeout configurable rather than fixed at 30 seconds would allow operators to tune the trade-off between detection speed and coordination overhead for their specific environment.

On the feature side, node selector rules would allow workloads to express hardware or topology requirements. For instance, a machine learning training job might require a GPU, and the scheduler should only consider agents that satisfy that constraint. Rolling workload updates, similar to what Kubernetes provides for deployments, would enable operators to update container images or configuration without downtime by progressively replacing old instances with new ones. A higher-level abstraction on top of individual workloads, analogous to ReplicaSets in Kubernetes, would let users declare a desired number of replicas and let the system handle creation, scaling, and reconciliation automatically. Finally, integrated telemetry, including metrics export, structured logging, and distributed tracing, would give operators visibility into scheduling decisions, agent health, and workload lifecycle events, which is essential for running any orchestration system in production.
